% !PRINTED 478
\weberpage{478}
Hat man irgend ein System nicht verschwindender Functionen $X_1, X_2, X_3$, so kann man daher solche Resolventen immer bilden. Durch das Formenproblem der Gruppe $G_{168}$ ist also dann die Gleichung 7\textsuperscript{ten} Grades mit der Gruppe $P_{168}$ zugleich gelöst.

Um die Lösung dieser Gleichungen 7\textsuperscript{ten} Grades auf das specielle Formenproblem, wie es bei den elliptischen Functionen auftritt (mit $f = 0$), zurückzuführen, ist dann noch eine accessorische Gleichung 4\textsuperscript{ten} Grades zu lösen, so wie wir, um die allgemeine Gleichung 5\textsuperscript{ten} Grades auf die Ikosaëdergleichung zurückzuführen, eine accessorische Quadratwurzel n\"othig fanden.

Alles kommt also jetzt noch darauf an, die Functionen $X_1, X_2, X_3$ der Wurzeln $\lambda$ diesen Forderungen gemäss zu bestimmen. Um dies zu ermöglichen, müssen wir einige einfache Sätze aus der allgemeinen Invariantentheorie benutzen, die wir im folgenden Paragraphen, soweit sie für unsere Aufgabe in Betracht kommen, ableiten wollen.

\begin{center}
§. 129.

\gesperrt{Contragrediente Gruppen.}
\end{center}

Wir haben schon im §.\,37 den Begriff der contragredienten Transformation erläutert. Sind nämlich

\weq{1}{A = \begin{pmatrix} a_1, & b_1, & c_1 \\ a_2, & b_2, & c_2 \\ a_3, & b_3, & c_3 \end{pmatrix}, \quad A_1 = \begin{pmatrix} a_1, & a_2, & a_3 \\ b_1, & b_2, & b_3 \\ c_1, & c_2, & c_3 \end{pmatrix}}

zwei zu einander transponirte Substitutionen, sind $x_1, x_2, x_3$ und $\xi_1, \xi_2, \xi_3$ zwei Reihen von Variablen, die durch die Substitutionen

\weq{2}{(y) = A(x), \quad \xi = A_1(\eta)}

in zwei neue Reihen von Variablen $y_1, y_2, y_3$ und $\eta_1, \eta_2, \eta_3$ transformirt werden, so haben wir diese beiden Reihen von Variablen und ebenso ihre Transformationen \gesperrt{contragredient} genannt.

Durch die Substitution $y = A(x)$ wird jede Function $\Phi(x_1, x_2, x_3)$ der $(x)$ in eine Function der $(y)$ transformirt, und die Bildung der Abgeleiteten ergiebt:

\weq{3}{\frac{\partial \Phi}{\partial x_1} = a_1 \frac{\partial \Phi}{\partial y_1} + a_2 \frac{\partial \Phi}{\partial y_2} + a_3 \frac{\partial \Phi}{\partial y_3}, \cdots}
%CONT>>
